%=============================================================================!
% E.1  THE ENVIRONMENT                                                        !
%=============================================================================!
\section{The environment}
\label{sec:env-environment}

Every calculation in the book runs in one Python environment, created with
the conda package manager from the file \texttt{environment.yml} at the root
of the course material. Conda installs compiled libraries as well as Python
packages, which matters here because the environment carries LAMMPS, the
PyTorch machine-learning library and the LaTeX-driven animation tool Manim,
none of which is pure Python. Two commands build it and install
\texttt{mdlab} into it:
\begin{code}[language=bash,numbers=none]
conda env create -f environment.yml
conda activate mdbook && pip install -e mdlab
\end{code}
The second command installs \texttt{mdlab} in editable form, so a change to
its source takes effect without reinstalling. The command \texttt{pytest
mdlab} then runs the tests of every routine built so far; they should all
pass before any notebook is opened.

The file names packages without versions and leaves the choice to the
solver, so two builds months apart can differ. \Cref{tab:env-versions}
records the versions against which the book's numbers were produced; the
script that prints it is in \texttt{scripts/appE\_environments}, and running
it in a new build shows at once what has moved.

\begin{table}[ht]
\centering
\caption{Versions of the principal packages in the environment that
produced the results of this book.}
\label{tab:env-versions}
\input{appendices/appE_environments/versions_table}
\end{table}

Some codes cannot share this environment. A package may require an exact
version of a library on which another package sets an incompatible lower
bound, and then no single environment satisfies both. The MACE potential
used from Chapter~12 onwards and the TrajCast propagator of Chapter~23 are
such a pair: MACE requires a particular early release of the
\texttt{e3nn} library for rotation-equivariant networks, and TrajCast a
later one. Each such code therefore receives its own environment, built in
the chapter that first runs it, and data pass between environments as
files.
